\documentclass[12pt]{article}
\usepackage{amsmath, amssymb, graphicx, listings, xcolor}
\usepackage[margin=1in]{geometry}
\usepackage{hypercap}

\definecolor{BLUE}{RGB}{0,84,159}
\definecolor{GRAY}{RGB}{128,128,128}

\begin{document}

\begin{center}
  {\Huge\bfseries\color{BLUE} ENEE322/ENEE323 Project 3}\\[6mm]
  {\Large Piano Note Spectral Analysis and Synthesis}\\[4mm]
  \hrule
  \vspace{5mm}
  {\large \textbf{Student:} Yahya Mouamil}\\[3mm]
  {\large \textbf{UID:} 119873526}\\[3mm]
  {\large \textbf{Date:} April 9, 2026}\\[3mm]
  \hrule
\end{center}

\section*{Problem (i): DFT and Fundamental Frequency}

The DFT of the piano note was computed. The magnitude spectrum revealed regularly spaced harmonic peaks.

\begin{figure}[h]
\centering
\includegraphics[width=0.7\textwidth]{Fig1_Spectrum.png}
\caption{Magnitude spectrum of the piano note}
\end{figure}

The fundamental frequency index was found to be:
\[
\boxed{K_1 = 1173}
\]

The corresponding frequency is:
\[
f_{\text{fundamental}} = \frac{K_1 \cdot f_s}{N} = \frac{1173 \cdot 11025}{33000} = 391.55 \text{ Hz}
\]

This corresponds to the musical note \boxed{\textbf{G4}}.

\section*{Problem (ii): dB Magnitude Plot}

The magnitude spectrum in decibels is shown below:

\begin{figure}[h]
\centering
\includegraphics[width=0.7\textwidth]{Fig2_dB.png}
\caption{dB magnitude spectrum}
\end{figure}

The maximum dB value is:
\[
\boxed{dB_{\text{max}} = 39.84 \text{ dB}}
\]

\section*{Problem (iii): dB Thresholding}

For attenuation values of 40 dB and 60 dB, entries below \(dB_{\text{max}} - att\) were zeroed out.

\begin{tabular}{|c|c|}
\hline
\textbf{att (dB)} & \textbf{Non-zero entries} \\
\hline
40 & 550 \\
60 & 3678 \\
\hline
\end{tabular}

The reconstructed signals without taper have abrupt boundaries causing audible clicks. Applying the taper smooths these discontinuities.

\section*{Problem (iv): Pure Harmonic Signal}

A pure harmonic signal was created using only frequencies at multiples of \(K_1 = 1173\):

\[
\text{Max harmonics possible} = 13
\]
\[
\text{Khalf} = 586
\]

\begin{figure}[h]
\centering
\includegraphics[width=0.7\textwidth]{Figure4_Vee2.png}
\caption{Vee2 - Pure harmonic signal with constant amplitude}
\end{figure}

\section*{Problem (v): Hamming Window}

Applying a Hamming window to Vee2 smooths the boundaries but does not recreate the natural decay envelope of the piano.

\section*{Problem (vi): Highest Harmonic Replacement}

Vsave1 contains \textbf{6 nontrivial harmonics}. The highest harmonic index is 7039. Replacing its cluster with a single entry produced a subtle but perceptible change in timbre.

\section*{Problem (vii): Exponential Envelope}

An exponential envelope \(A e^{-ct/N}\) with \(A = \max(|v_{\text{note}}|)\) and \(c = 3.5\) was fitted to the original signal:

\begin{figure}[h]
\centering
\includegraphics[width=0.7\textwidth]{Figure7_Envelope.png}
\caption{Exponential envelope fit to original vnote}
\end{figure}

Multiplying Vee2 by this envelope produced Vee4, which sounds significantly closer to the original piano note, though it lacks the inharmonicity and noise of a real piano.

\begin{figure}[h]
\centering
\includegraphics[width=0.7\textwidth]{Figure8_Vee4_Spectrum.png}
\caption{Magnitude spectrum of Vee4}
\end{figure}

\section*{Summary of Results}

\begin{tabular}{|l|l|}
\hline
\textbf{Parameter} & \textbf{Value} \\
\hline
\(K_1\) & 1173 \\
Fundamental frequency & 391.55 Hz \\
Musical note & G4 \\
\(dB_{\text{max}}\) & 39.84 dB \\
Harmonics in Vsave1 & 6 \\
Highest harmonic index & 7039 \\
Max harmonics possible & 13 \\
\hline
\end{tabular}

\vspace{8mm}
\begin{center}
  {\color{BLUE} \Large End of Project 3}
\end{center}

\end{document}
